\documentclass{article}

\usepackage[preprint]{neurips_2026}

\usepackage[utf8]{inputenc}
\usepackage[T1]{fontenc}
\usepackage{hyperref}
\usepackage{url}
\usepackage{booktabs}
\usepackage{amsfonts}
\usepackage{amsmath}
\usepackage{amssymb}
\usepackage{amsthm}
\usepackage{microtype}
\usepackage{xcolor}
\usepackage{graphicx}

\newtheorem{definition}{Definition}
\newtheorem{assumption}{Assumption}
\newtheorem{theorem}{Theorem}
\newtheorem{proposition}{Proposition}

\title{Proper Time for Verified Model Selection}

\author{%
  Anonymous Authors \\
  \texttt{}
}

\begin{document}

\maketitle

\begin{abstract}
\end{abstract}

\section{Introduction}

Achille and Soatto argue that agents should be understood as task solvers whose
performance is measured by \emph{time to solution}: the resource required to
produce a verified artifact, not the model's self-assessment of progress
\citep{achille2025agents}.  This perspective changes a basic practical
question.  Modern model choice is usually treated as a leaderboard problem: run
each candidate on a benchmark and pick the one with the highest average score.
For verified tasks, this is often the wrong unit of comparison.  In code
generation, repository repair, formal reasoning, and exact-answer reasoning
benchmarks, a system is useful only when it produces an artifact accepted by an
external verifier.  The deployment question is therefore not merely whether a
model can eventually solve an instance, but how much wall-clock time, token
budget, and sampling effort must be spent before a verified solution appears.

This distinction matters because verified performance is trajectory-valued.  A
small model may be cheap enough that several retries are preferable to one call
to a larger model.  A larger model may be slower but decisive on modes where
smaller models repeatedly fail.  A router may observe enough information about a
fresh instance to select the right worker, or it may waste that information by
choosing an action that is not on the mode-conditioned frontier.  Average
accuracy collapses these mechanisms into one number, even though they imply
different engineering decisions.

The scientific reason to introduce task modes is the same reason algorithm
selection and portfolio methods exist: instance structure can change which
solver is cost-effective.  In Rice's algorithm-selection formulation, one
chooses an algorithm using observable features of the problem instance
\citep{rice1976algorithm}; SATzilla and AutoFolio made this idea practical by
learning when different solvers dominate on different regions of a task
distribution \citep{xu2012satzilla,lindauer2015autofolio}.  In machine
learning, transferability and task-embedding work similarly treats tasks as
structured objects rather than interchangeable samples
\citep{achille2019task2vec}.  Our modes play this role for verified language
model deployment.  They are not post-hoc names for success or failure.  They
are predeclared, solver-relevant partitions or covariates---for example
difficulty strata, algorithmic categories, or reasoning families---under which
different models may have different cost--success curves.

This makes a concrete prediction.  If one model has the best cost-adjusted
first-passage profile in every mode, routing is trivial: no mode label can make
a different model worth running.  If the frontier crosses across modes, then
model choice becomes conditional.  A smaller or cheaper model can be optimal
under a wall-clock, token, or price budget on modes where its success rate is
high enough; a larger model can be optimal on modes where the smaller model's
retry curve saturates below the required reliability.  This is the verified
model-selection analogue of comparative advantage.

We do not introduce proper time as a new philosophical object.  Instead, we
make the Achille--Soatto principle empirical.  Given a fixed verifier, a
resource clock, and repeated stochastic runs, we estimate the first-passage
curve
\[
    F_M(t) = \Pr[\text{model } M \text{ has produced a verified solution by resource } t],
\]
keeping unsolved instances as right-censored observations.  From this curve we
study the operating frontier
\[
    \tau^\star_M = \inf_t \frac{t}{F_M(t)},
\]
which asks how much resource must be paid per unit of verified success at the
best operating point.  This turns model selection into verified first-passage
accounting rather than benchmark racing.

The main observation is that proper-time differences are not monolithic.  A
system can improve because its attempts are cheaper, because its worker has
higher verified competence on a task mode, because the instance exposes
information about which mode it belongs to, or because a router allocates that
information well.  We therefore organize verified model selection around four
diagnostic quantities:
\begin{enumerate}
  \item \textbf{Cost}: the resource consumed by an attempt or by a
  stop-at-first-success trajectory.
  \item \textbf{Competence}: the mode-conditioned first-passage success curve of
  a worker under a fixed verifier.
  \item \textbf{Information}: the value of knowing a mode or router-visible
  signal before deployment.
  \item \textbf{Allocation regret}: the loss of the implemented router relative
  to the best action available under the same information.
\end{enumerate}
The resulting decomposition separates two questions that are often conflated:
whether there is useful specialization in the model menu, and whether a router
actually exploits it.

We instantiate this protocol on verified coding and reasoning tasks.  For each
worker model we log every sampling attempt, the per-attempt resource usage, the
verifier outcome, and the cumulative resource to the first passing attempt; runs
that never pass remain in the data as censored observations.  This trace format
lets us evaluate single-model frontiers, retry policies, oracle allocations,
statistical oracles that learn from held-out folds, and prompt-only LLM routers
without rerunning workers.  The router never writes a solution; it only chooses
an allocation over already logged worker attempts, and is judged by the verified
time-to-solution that allocation would have incurred.

Our contributions are:
\begin{enumerate}
  \item We operationalize proper time for verified model selection as an
  estimable first-passage object with wall-clock and token clocks.
  \item We give a diagnostic decomposition of proper-time gains into cost,
  competence, information value, and allocation regret.
  \item We test the retry-crossover prediction empirically, distinguishing the
  closed-form Bernoulli approximation from the heterogeneous first-passage
  curves observed in real task sets.
  \item We evaluate when routing is non-trivial: a router can beat the best
  single model only when task modes induce cost-adjusted frontier crossings, and
  only if the router uses the available information well.
\end{enumerate}

The paper's central message is simple: accuracy tells us whether a model can
solve; proper time tells us whether it is worth running.

\section{Verified Tasks and First Passage}
\label{sec:formalism}

We begin with the object on which proper time is measured.  The task must have
an external success condition: the model may propose artifacts, but the
acceptance event is determined by a verifier outside the model.

\begin{definition}[Verified task]
\label{def:verified-task}
A verified task is a tuple
\[
    \mathcal T = (\mathcal X,\mathcal Y,\mu,V),
\]
where \(\mathcal X\) is an instance space, \(\mathcal Y\) is an artifact space,
\(\mu\) is a distribution over instances, and
\[
    V:\mathcal X\times\mathcal Y\to\{0,1\}
\]
is an external verifier.  For an instance \(x\in\mathcal X\) and candidate
artifact \(y\in\mathcal Y\), the event \(V(x,y)=1\) means that \(y\) is accepted
as a verified solution to \(x\).
\end{definition}

The definition is intentionally agnostic about the kind of artifact.  In a code
task, \(x\) is a programming problem and \(y\) is a generated program; the
verifier runs the declared tests.  In an exact-answer reasoning task, \(x\) is a
question and \(y\) is the model's final answer; the verifier applies a frozen
extraction and normalization rule and checks exact match.  Scalar evaluators can
be included by thresholding: if \(R(x,y)\in\mathbb R\) is an external reward or
loss, then \(V_\gamma(x,y)=\mathbf 1\{R(x,y)\ge\gamma\}\) for a predeclared
threshold \(\gamma\).

Two points are important.  First, verification is external: success is not the
model's claim that the answer is correct.  Second, the verifier defines the
unit of progress.  Any comparison of models, retries, or routers in this paper
is made relative to the same task distribution and the same frozen verifier.

\begin{definition}[Deployment action and resource clock]
\label{def:action-clock}
Let \(\mathcal A\) be a finite set of deployment actions.  An action
\(a\in\mathcal A\) is any fixed procedure that, given an instance \(x\) and
internal randomness \(r\in\Omega\), produces candidate artifacts over time.
Examples include a single language-model call, a language model wrapped in a
harness, a repeated-sampling policy, a sequential allocation across models, a
multi-agent system, or a broader orchestration capable of solving the instance.

A resource clock is a nondecreasing map
\[
    C^a_k(x,r)\in\mathbb R_{\ge 0}, \qquad k=1,2,\ldots,
\]
where \(C^a_k(x,r)\) is the cumulative resource spent by action \(a\) after its
\(k\)-th candidate artifact on instance \(x\).  The resource may be wall-clock
seconds, generated tokens, dollar cost, verifier calls, or a fixed weighted
combination of these quantities.  We assume \(0<C^a_1\le C^a_2\le\cdots\).
\end{definition}

\begin{definition}[Resource-indexed trajectory]
\label{def:resource-trajectory}
For \(a\in\mathcal A\), instance \(x\), and randomness \(r\), write
\[
    \bigl(C^a_k(x,r),Y^a_k(x,r)\bigr)_{k\ge 1}
\]
for the sequence of cumulative resources and candidate artifacts produced by
the action.  The associated resource-indexed artifact process is
\[
    Y^a_c(x,r)
    :=
    Y^a_{K_a(c;x,r)}(x,r),
    \qquad
    K_a(c;x,r)
    :=
    \max\{k:C^a_k(x,r)\le c\},
\]
with \(Y^a_c\) undefined when the set is empty.  Thus \(Y^a_c\) is the latest
artifact available after spending resource \(c\).
\end{definition}

The cumulative convention is part of the definition.  If a repeated-sampling
action succeeds on its third candidate, the cost of the first two failed
candidates is still paid.  Likewise, if a sequential allocation tries a cheap
worker before an expensive one, the verified time of the allocation includes
every consumed attempt up to and including the first accepted artifact.

\bibliographystyle{unsrtnat}
\bibliography{references}

\end{document}
